\section{Présentation générale du SHOM}

Le Service hydrographique et océanographique de la Marine, plus couramment appelé SHOM, est l'organisme français de référence dans les domaines de l'hydrographie, de l'océanographie et de l'information géographique maritime et littorale. Il s'agit d'un établissement public à caractère administratif placé sous la tutelle du ministère des Armées.

Le SHOM trouve ses origines dans le Dépôt des cartes et plans de la Marine, créé en 1720. Cette ancienneté traduit l'importance historique de la connaissance du milieu marin pour la navigation, la sécurité maritime et les activités de Défense. Son rôle a progressivement évolué avec les besoins de la navigation, les avancées technologiques, la numérisation des données et le développement de nouveaux usages liés à la mer et au littoral.

Aujourd'hui, le SHOM exerce trois grandes missions. La première est l'hydrographie nationale, qui consiste à produire et maintenir les informations nécessaires à la sécurité de la navigation maritime. La deuxième concerne le soutien hydro-océanographique à la Défense, afin de fournir aux forces armées françaises des informations utiles à la préparation et à la conduite des opérations navales. La troisième porte sur le soutien aux politiques publiques de la mer et du littoral, notamment à travers la mise à disposition de données et d'expertises pour la gestion des espaces maritimes, la prévention des risques et la protection de l'environnement marin.

\section{Missions, produits et données}

Les activités du SHOM reposent sur la collecte, le traitement, la validation, la production et la diffusion de données maritimes. Ces données peuvent concerner la bathymétrie, les marées, les courants, la nature des fonds, les épaves, les limites maritimes ou encore les caractéristiques physiques de l'océan. Elles sont utilisées par des navigateurs, des institutions publiques, des collectivités, des chercheurs, des acteurs économiques et les forces armées.

La cartographie marine constitue une production centrale du SHOM. À partir des données collectées et validées, l'établissement produit des cartes marines papier ainsi que des cartes électroniques de navigation, ou ENC, conformes aux normes internationales de l'Organisation hydrographique internationale. Ces cartes représentent les côtes, les profondeurs, les dangers, les aides à la navigation et les informations nécessaires à la sécurité des navigateurs.

Le SHOM produit également des ouvrages nautiques complémentaires aux cartes marines. Les instructions nautiques décrivent les côtes, les ports, les dangers, les conditions locales de navigation et les zones réglementées. Le livre des feux répertorie les feux et signaux lumineux utiles à la navigation. Les radiosignaux regroupent quant à eux les informations radioélectriques nécessaires à la sécurité et à la communication maritime.

Ces productions existent aujourd'hui dans un contexte de transition numérique. Les cartes électroniques, les services web, les portails de diffusion et les bases de données occupent une place croissante, mais les cartes papier et certains ouvrages traditionnels restent encore utilisés dans plusieurs contextes. Cette coexistence entre productions historiques et numériques crée des enjeux importants d'harmonisation, de qualité, de traçabilité et de maintenance des chaînes de production.

\section{Environnement numérique et contraintes techniques}

En raison de son rattachement au ministère des Armées et de la nature des données traitées, le SHOM évolue dans un environnement numérique fortement contraint. Les enjeux de cybersécurité, de souveraineté des données, de disponibilité des systèmes et de maîtrise des chaînes de production y occupent une place centrale.

Une part importante de l'infrastructure informatique est hébergée et maintenue en interne. Les serveurs, bases de données, applications métiers et partages de fichiers sont exploités dans un environnement contrôlé. Les droits d'accès sont définis selon les utilisateurs, les groupes, les services et les responsabilités associées. Certains environnements sont dédiés à la production, tandis que d'autres permettent les développements, les tests ou les manipulations techniques.

Les développements réalisés dans ce contexte ne peuvent donc pas être pensés comme des applications isolées. Ils doivent s'intégrer dans un système d'information existant, composé de chaînes historiques, de dépendances fortes avec les métiers, de données volumineuses, de formats spécialisés et de contraintes de sécurité. Cette réalité influence directement les choix techniques, les méthodes de développement et la manière de faire évoluer les outils internes.

\section{Service d'accueil et place dans l'organisation}

À mon arrivée au SHOM, mon alternance se déroulait au sein du pôle SYSPROD, rattaché à la division Produits et services maritimes de la Direction des opérations, de la production et des services. Ce pôle était chargé de la maintenance, de l'expertise et de l'évolution des systèmes de production utilisés dans les activités nautiques et cartographiques du SHOM.

Les missions du pôle concernaient notamment l'étude de besoins, le développement de maquettes, le déploiement et l'interconnexion de briques applicatives, les services web, la documentation technique et l'amélioration des chaînes existantes. Ces activités s'inscrivaient dans un contexte où les outils informatiques devaient répondre à des contraintes métier fortes, tout en restant maintenables, fiables et adaptés aux processus de production.

Environ un an après mon arrivée, le pôle SYSPROD a été dissous dans le cadre d'une réorganisation interne. Le pôle DATA a alors été créé, en regroupant l'ensemble des membres de SYSPROD ainsi que plusieurs agents issus d'autres pôles. Cette évolution organisationnelle traduit l'importance croissante des problématiques liées à la donnée, à sa structuration, à sa valorisation et à l'amélioration des chaînes de traitement au sein du SHOM.

J'ai ainsi poursuivi mon alternance au sein du pôle DATA, dans une équipe aux profils variés, impliquée dans la création, la maintenance et l'évolution des chaînes de production utilisées pour les ouvrages nautiques et cartographiques. Ma position dans cette équipe m'a permis de participer à des travaux d'analyse, de développement, de documentation et de modernisation progressive de traitements existants.

Cette place dans l'organisation est importante pour la suite du dossier, car mes missions ne relèvent pas uniquement du développement logiciel au sens strict. Elles impliquent également la compréhension de processus métier complexes, l'analyse de systèmes existants, la prise en compte des contraintes de production, la sécurisation des traitements, l'amélioration de la qualité logicielle et la transmission de connaissances techniques.